\include{zLaTex-Preamble}

\title{Concurrent Hash Tables - CMSC 23000 } 
\author{Joshua Miller}
%% \date{}

\begin{document}
\rhead{\fancyplain{}{Josh Miller}}
\maketitle


\renewcommand{\t}[1]{{\tt\sc #1}}

\section{General Design}

{\bf Modules}
\begin{itemize}
\item util.h -- contains macros and definitions for error handling and debugging
\item hasht.h, libhasht.a -- contains definitions for each of the four hash table types 
\item tests.c -- contains the tests for the entire package
\item app.h -- contains the different hash performance tests
  \t{Serial, Parallel, ParallelNoLoad}.
\item experiment\{1,2,3,4\}.c -- contains the code to run the experiments.
\end{itemize}

\section{General}

The first comment I would like to make about this project is that the
rates I gathered during the experiments varied an extereme amount.  

%% INSERT STANDARD DEVIATION ANALYSIS OF EXPERIMENT 1

The hash table created for this project are contained in the
statically linked library {\tt libhasht.a} located in {\tt src/libhasht} with
{\tt hasht} being the prefix for all cusom hash tables.  The general
design of the hash tables if a polymorphic inheritance from the main
struct \t{hasht\_t}.  This 'object' contains pointers to it's
initialization, add, remove, resize, and contains functions.  Included
below is the majority of \t{hasht.h}.  I apologize for the inclusion a
nearly complete header file, but I feel that the information included
in it is an encompassing description of the \t{hasht}
abstraction. Each 'subclass' is described in later sections.


\lang{C}
\begin{code}
/* The top level hash table abstraction */
typedef struct hasht_t hasht_t;   

/* Pointer types for each table interaction member */
typedef int (*hasht_init_f)    (hasht_t *, int, int);
typedef int (*hasht_free_f)    (hasht_t *);
typedef int (*hasht_get_type_f)(hasht_t *);
typedef int (*hasht_add_f)     (hasht_t *, void*, int);
typedef int (*hasht_contains_f)(hasht_t *, int);
typedef void *(*hasht_remove_f)  (hasht_t *, int);
typedef void *(*hasht_resize_f)  (hasht_t *);

/* Enumeration to be contained in hasht object */
typedef enum hasht_type_t 
{
    LOCKING,
    LOCKFREEC,
    LINEAR,
    AWESOME
} hasht_type_t;

typedef struct hasht_locking_t
{
    ll_t **buckets;
    pthread_rwlock_t *rwlocks;
    int n_rwlocks;
} hasht_locking_t;

typedef struct hasht_lockfreec_t
{
    ll_t **buckets;
    pthread_rwlock_t *rwlocks;
    int n_rwlocks;
} hasht_lockfreec_t;

typedef struct linear_node_t{
    unsigned int key;
    unsigned char inuse;
    void *data;
    unsigned int steps;
} linear_node_t;

typedef struct hasht_linear_t
{
    int n_locks;
    int resizing;
    pthread_mutex_t *locks;
    linear_node_t *buckets;
} hasht_linear_t;

typedef struct awesome_node_t awesome_node_t;

struct awesome_node_t{
    unsigned int residue;
    unsigned int key;
    void *data;
    awesome_node_t *child;
    unsigned char inuse;
    unsigned char valid;
};

typedef struct hasht_awesome_t
{
    awesome_node_t root;
 } hasht_awesome_t;

struct hasht_t 
{
    /* Type specifier */
    hasht_type_t type;

    /* Scalar properties */
    int logsize;
    int capacity;
    int count;

    /* Member Function Pointers */
    hasht_add_f      add;
    hasht_contains_f contains;
    hasht_remove_f   remove;
    hasht_resize_f   resize;

    /* Member types */
    hasht_locking_t      locking;
    hasht_lockfreec_t    lockfreec;
    hasht_linear_t       linear;
    hasht_awesome_t      awesome;
}; 


/* Constructor for new hash table */
hasht_t *hasht_new(hasht_type_t type, int capacity, int expected_threads);
int hasht_locking_init(hasht_t *table, int capacity, int expected_threads);

\end{code}

To avoid the cost of freeing memory during computation. I wrote a
rudimentary garbage collector. During computation, allocations are
stored in a reference array.  After computation a call to
\t{garbageCollect()} frees all of the memory allocated by the hash
tables.


\section{Lock-based Closed-address}

The \t{hasht\_locking\_t} table contains the following
\begin{itemize}
\item all inherited properties from \t{hasht\_t}
\item an array of linked lists of size \t{MAX\_DEPTH}
\item an array of read write locks of size \t{nlocks}
\item an integer \t{Count} specifying how many packets are in the table
\end{itemize}

Because we can stripe the hash table, the locks are held in parallel
array, such that the lock for item $i$ is accessed by calculating
$i \mod$\t{nlocks}.

Editing (writing, resizing) the hash table requires the effective
acquisition of both the write and read locks. Reading (contains) from
the hash table requires acquisition of only the read locks.  This
means that for reads and writes will interfere with each other, this
being the largest source of contention and overhead.  The difficulties
in implementing this were slightly more challenging than implementing
the \t{LOCKFREEC} because the later consisted more of modifications
than original design. The resize on the \t{LOCKING} table is
accomplished by acquiring all the locks, creating a new and duplicate
bucket array, then redirecting the table's bucket array to the newly
created one.  This is all sorts of painful because (1) we have to lock
the entire table and (2) we have to go through the trouble of
re-adding each element to the new bucket array.  Avoiding these two
factors would eventually become part of the design requirement for my
\t{AWESOME} table.  Furthermore, it was important to verify that the
table had not been resized while acquiring the write lock for an add
or remove call. If it had been, then the respective write call would
simply be recusrively called again after freeing any obtained locks. A
resize call that acquired the locks during an second extern resize
would simply return.

\section{Lock-Free-Contains Closed-Address}


The \t{hasht\_lockfreec\_t} table contains the following
\begin{itemize}
\item all inherited properties from \t{hasht\_t}
\item an array of linked lists of size \t{MAX\_DEPTH}
\item an array of \t{pthread\_mutex\_t} of size \t{nlocks}
\item an integer \t{Count} specifying how many packets are in the table
\end{itemize}

Again, the locks are held in parallel array, such that the lock for
item $i$ is accessed by calculating $i \mod$\t{nlocks}.

The major design changes from the \t{hasht\_locking\_t} were (1)
checking the validity of the table before returning from a contains
call and (2) making sure that nodes were not deleted while a contains
was attempinting to traverse a list.  The later objective was easily
obtained because my implementation of the locking table lent itself an
atomic replacement of the entire bucket list.  The former objective
was acheived simply by maintaining a record of the table size at the
beginning of the call to contains and comparing before exiting.  If
this test fails, then the contains method is simply called again.

\section{Linearly Probed Open-Address}

The table contains the following
\begin{itemize}
\item a single array for nodes each node containing
\subitem  - a pointer to the data
\subitem  - a key
\subitem  - a binary setting specifying availability
\subitem  - a step count
\item a size
\end{itemize}

Finding an element in the linear probe hash table involved masking the
key, and selecting the respecitve node from the array.  After a call
to \t{add()} acquiring the lock (again a striped locking scheme) if
the node is available then modify the availibility flag to signify
that it is now taken. If we have to search along the array, then we
note in the original location how far we had to travel, or the maximum
if another call has already noted its travel distance. I made the
silly mistake of setting this counter to 0 on a \t{remove()} call,
which caused \t{contains()} to return incorrectly.

When a resize happens, the locks for the entire array are acquired and
a new array of nodes are created. The old items are then re-added to
the new array before the locks are released.

\section{The most awesomest hash table that concurrent
  data structures has ever seen, a.k.a AWESOME}


The two main focal points that I wanted to improve were (1) locking
and (2) resizing (3) locality.  All of the above designs required a
stop-the-world resizing strategy, and I wanted to create a design that
was free of this long pause in progress.  Furthermore, I wanted to
maximize the spacial locality of my data structure; I decided to use
arrays, rather than linked lists, as the backbone for this hash
table. The problem with arrays, however, as we saw with the linear
probe is the resize. For these reasons I came up with the following
design. The idea is a dynamicly multidimensional hash table somewhere
between a linear probe and a tree.

Consider the table to be one dimensional at initialization. Add calls
proceed as if the table were a linear probe hash table, except each
call is only allowed to walk down the array a number of times
specified by the \t{OPTIMISM} parameter. This parameter specifies how
optimisitic there will be space in the current dimension.  If an open
position is found within this limit, the item is inserted as in a
linear probe.  However, if there is no longer room within the limit of
the initial hash position, a new dimension is added.  This dimension
branches off of the node itself. In this way it is like a tree, but
each branch is not a single item but a 'dimension' with a preset
length \t{DIM\_LEN}. Therefore, each dimension consists of an array of
nodes, which may or may not have dimensions branching off of them.  If
the algorithm is unable to find a desired location/item in the current
dimension, then the algorithm steps into the next dimension.  Because
each item in a new dimension will have the same index in the old
dimension, a new array index is calculated by bit shifting the key to
the right by the dimension level and remasking on \t{DIM\_LEN}.  

Through the use of atomic operations and recursively incremental
resizes, I was able to make the table completely lock free. There is
no resize call, an add call simply inserts a dimension if necessary.
N.B: The new dimension must be attached with an atomic compare and set
call, which prevents replacing a dimension allocated and attached by a
concurrent \t{Add()} call. 

The trade off of this table is (a) locality, a lock-free nature, and
quicker tree-like lookup versus (b) pre-allocating node arrays which
may be in excess. Therefore I ran an extra experiment to optimize the
parameters \t{OPTIMISM} and \t{DIM\_LEN}. A plot of the results are
presented below.  The results from this experiment were used as
settings for the \t{Awesome} lock in successive experiments.

Moreover, because of the tree-like nature and the shifted address
masking, if any given key is contained in the table, we know exactly
which nodes it had to pass through at each dimension of the
table. Therefore, I kept a \t{residue} on each node.  As an \t{Add()}
call either passed through a node to a new dimension, or assigned an
item to the node, it atomically \t{or}-ed the key onto the residue.
In this way, a contains or remove call knows in advance if a key is
absolutely \emph{not} present in recursive levels of the table. This
however, does not help a lookup if the key has been added and
successively removed from the table.


\section{Tests}

Each test follows the same testing scheme through progressively more
rigorous composite tests.  First the add and remove calls are tested
by verifying containment for small numbers of elements.  Next, a
battery of random keys are added to the table. Then a similar scheme
but with more than one thread. The final test walks over a key space
in a pseudo-random cycle 4 times.  Each key is added, with a random
possibility of being removed afterwards.  The testing function keeps
track of and tests against containment in an array parallel to the key
array. I am very satisfied with this test suite. 

To run the tests
\lang{bash}
\begin{code}
> make tests
> ./test
\end{code}

Or to test a specific table
\begin{code}
> ./test LOCKING
> ./test LOCKFREEC
> ./test LINEAR
> ./test AWESOME
\end{code}

Expected output:
\begin{code}
TEST: [                 add/remove] [ serial_simple_hash_test(AWESOME)] [ PASSED ]
TEST: [          add/remove/resize] [ serial_resize_hash_test(AWESOME)] [ PASSED ]
TEST: [         and/resize/contain] [     parallel_hash_test1(AWESOME)] [ PASSED ]
TEST: [  random traversal of key..] [     parallel_hash_test2(AWESOME)] [ PASSED ]
\end{code}


\section{Experiments}

\subsection{Dispatcher Rate}

This experiment involves running the parallel test with no load and
minimal packet size for the number of threads equal to the number of
cores. This will present us with the rate at which the dispatcher
thread can enqueue the packets onto the lamport queues. I predicted
that the dispatcher rate will be based upon the relative performanc of
the \t{hashgenerator}, which I believe will be slightly slower, and
therefore the performance will fall just below the performance of the
dispatcher rate in previous projects. This prediction was true.

I decided to use the \t{HomeQueue} work load method because each
thread should have a constant work load supplied to it, as the
dispatch dishes one packet out to each queue iteratively. Furthermore,
since this is a time based test, the queue fairness is less important
in the performance statistics.

When running this test, I decided that I was looking for the quickest
possible rate that the dispatcher could work, otherwise I could get a
hash table performance that is better than the 'ideal' performance
that this experiment is simulating.  Therefore I chose the maximum
dispatcher rate at each of the thread counts I would run Experiment~3
at.  This allowed me to compare each data point to an 'ideal' speedup.

\figH {../plots/plot1.png}{Dispatcher rate as a function of thread count}{plot1}{.9}

\subsection{Parallel Overhead}

Configuration: $n = 1$, $W = 4000$ 

This experiment involved running the parallel test with each of the
parallel hash tables, but on a single thread.  In comparison with the
performance of the serial hash table the speedup will be less than
1. 

\emph{Hypothesis: } I predicted that for all hash tables, the speedup
will be lower for the heavy write runs.  This is because the writing
has an added complexity relative to the serial version to handle
concurrency that the reading does not.

\shadetable{|c|c|c|}{Table Type & Rate & Speedup}
{Mostly Reads}{  
\t{Serial} & - & - \\
\t{Locking} & 59.57    &   \\
\t{LockFreeC} & 51.47  &   \\
\t{Linear} & 57.99     &   \\
\t{Awesome} & 53.92    &   
}{tab:test1}

\shadetable{|c|c|c|}{Table Type & Rate & Speedup}
{Heavy Writes}{  
\t{Serial} & - & - \\
\t{Locking} &  55.64   &   \\
\t{LockFreeC} & 80.21   &   \\
\t{Linear} &   51.85   &   \\
\t{Awesome} & 52.09     &   
}{tab:test1}

\emph{Result: }


\subsection{Speedup}

Configuration: $n \in \{1, 2, 4, 8, 16, 32\}$, $W = 4000$ with varying load schemes.

This experiment involved running a scaling test on all the parallel
hash packet implementations.  

\emph{Hypothesis: } I predicted that our program will be able
to acheive 80\% of ideal speedup, or just over 12 times the serial
version on 16 cores.  

\subsection{Mostly Reads}

\figH {../plots/plot3_0.png}{Table speedup comparision $P_+ = .09$ $P_- = .01$ $\rho = .5$ }
      {plot3.0}{.9}

\figH {../plots/plot3_1.png}{Table speedup comparision $P_+ = .09$ $P_- = .01$ $\rho = .75$ }
      {plot3.0}{.9}

\figH {../plots/plot3_2.png}{Table speedup comparision $P_+ = .09$ $P_- = .01$ $\rho = .9$ }
      {plot3.0}{.9}

\figH {../plots/plot3_3.png}{Table speedup comparision $P_+ = .09$ $P_- = .01$ $\rho = .99$ }
      {plot3.0}{.9}



\subsection{Heavy Writes}




\end{document}
